\label{sec:integer12}

The preceding positive-domain constructions and the following
integer-domain constructions are separate points in the comparison.
Here all twelve unknowns range over $\Z$.  We give two polynomial
coordinate changes of the pinned repository's current scalar-carrier
equations.  The transformations preserve the entire integer zero sets
bijectively and improve their degree--cost tradeoffs.  Centering is already
used elsewhere in the repository; the contribution here is its application
to the present scalar four-root carrier, with complete circuits and attained
degree certificates.

\subsection{Fixed coefficients and complete equations}

The ordered unknowns are
\[
 (j,C,z,h,r,v,w,T,c,n,t,s)\in\Z^{12}.
\]
Use the coherent \texttt{mixed-parity-unrounded-1064921} tuple, independently
of the compact44 tuple used by the positive-domain constructions:
\[
 \beta=2^{1064921}-1,\qquad
 \kappa_1=2^{1113},\qquad \kappa_2=2^{666104}.
\]
The remaining fixed integers $K_0,M_C,M_F$ are the exact source34 values
in the included coefficient payloads.  The complete finite recipe,
literal machine seed, and all 943 tiles are included; the verifier rebuilds
every coefficient and compares all bytes.  In particular $M_F$ is already
shifted by $\beta$.  The six coefficients are never mixed with another
tuple.

Define the following polynomial abbreviations; none is an additional
unknown:
\begin{align*}
 J&=6j+1,&q&=\beta J+1,&W&=12w+4,\\
 X&=qW,&Y&=(2v+1)q^3,&E&=XY,\\
 a&=E+Y,&H&=4a+3,&\Delta&=a^2+H,\\
 F&=(K_0+W)C-z(q-1),&&&
\end{align*}
\begin{align*}
 R&=\bigl(q(q-F)-C\bigr)(q^2-1)+(M_C+qM_F)J,\\
 k&=R+1+hE,\qquad \eta=c-Yk,\qquad \zeta=k-\eta,\\
 G&=(\kappa_1F)(\kappa_2C)-(q-\kappa_1F-\kappa_2C)^2,\\
 \Pi&=G\eta\zeta-(t^2-G-\eta\zeta+1)^2,\\
 D_g&=(3\Pi-4)s^2+1,\qquad I=src^2,\\
 D_f&=1+\Delta I^2,\qquad Q=\Delta^2I^2,\qquad\delta=Q-1,\\
 V&=D_f(c-R+4IT)-c,\qquad C_y=QV^2-1,\\
 \ell&=XY^2,\qquad D_t=1+\ell(\ell+1)k^2,\\
 K&=X(X+2ac)-1,\qquad S=IH^5,\qquad z_c=HT+K,\\
 A_c&=(z_c^2+3H^2)(z_c^3-19H^3).
\end{align*}
Set
\[
 r_a=9S^2D_f,\qquad r_b=81S^2D_t,\qquad r_e=729S^2D_g.
\]
Let $\mathcal F_{16}(Z;u_1,u_2,u_3,u_4)$ be the product of
$Z-\sum_{i=1}^4\epsilon_i\sqrt{u_i}$ over the sixteen choices
$\epsilon\in\{-1,1\}^4$.  Sign invariance makes this an integer
polynomial in its five arguments; every radicand has exponent at most
eight.  Our two equations are
\begin{align}
 P_R&=\delta^8\mathcal F_{16}
 \left(S(n+V)-A_c; r_a,r_b,r_e,\frac{S^2C_y}{\delta}\right)=0,
 \label{eq:s12-rational}\\
 P_I&=\mathcal F_{16}
 \left(S(n+\delta V)-A_c; r_a,r_b,r_e,S^2\delta C_y\right)=0.
 \label{eq:s12-integral}
\end{align}
The first formula means its symbolic polynomial after denominator
clearing, not numerical evaluation by division.  Both literal JSON
circuits contain only fixed integers, addition, subtraction, and
multiplication.

For completeness, the following polynomial recursion removes every
radical from the notation.  For a center $Z$, put $d=\delta$ and
$c_n=S^2C_y$ in~\eqref{eq:s12-rational}, or $d=1$ and
$c_n=S^2\delta C_y$ in~\eqref{eq:s12-integral}.  Define
\begin{align*}
 t_0&=Z^2+r_e,& t_1&=-2Z,& b_0&=r_a+r_b,\\
 q_0&=d(t_0-b_0)-c_n,&q_1&=dt_1,\\
 T_0&=d^2r_ar_b,&R_0&=dc_n,\\
 v_0&=q_0^2+r_eq_1^2-4(T_0+R_0b_0),&v_1&=2q_0q_1,\\
 z_{64}&=64T_0R_0,&&\\
 n_0&=v_0^2+r_ev_1^2-z_{64}t_0,&
 n_1&=2v_0v_1-z_{64}t_1.
\end{align*}
Then the polynomial output is $n_0^2-r_en_1^2$.  The independent oracle
checks all 495 coefficients against the sixteen-conjugate product.

\subsection{Unconditional integer zero-set bijections}

Write $P_R^{\rm old}$ and $P_I^{\rm old}$ for the pinned parent equations
with the same prefix and center $Sn-A_c$.  Every other abbreviation is
independent of $n$.  The maps
\[
 T_R:\ n_{\rm old}=n+V,\qquad
 T_I:\ n_{\rm old}=n+\delta V
\]
fix eleven coordinates.  Their inverse maps subtract $V$ and $\delta V$,
respectively.  These are mutually inverse polynomial maps of $\Z^{12}$,
without sign, nonzero, divisibility, or zero-equation hypotheses.
Direct substitution gives the identities
\[
 P_R=P_R^{\rm old}\circ T_R,\qquad
 P_I=P_I^{\rm old}\circ T_I.
\]
Consequently each pair has bijective complete integer zero sets, including
all exceptional branches as formal polynomials.  Composing these elementary
identities with the unchanged ordinary parent calibration gives, in the
same formalizability sense,
\[
 \mathrm{ZFC}\vdash
 \bigl(\forall x\in\Z^{12}\ P_R(x)\ne0\bigr)
 \leftrightarrow\operatorname{Con}(\mathrm{ZFC}),
\]
and the identical statement with $P_I$.  This composition introduces no
new native compiler claim.  The retained parent proof and its dependencies
remain part of the evidence; this section does not claim a new
machine-checked object-ZFC derivation.

\subsection{Exact degree certificates}

Only quantified coordinates contribute to total degree.  The prefix gives
\[
 \deg S=34,\quad\deg V=25,\quad\deg\delta=32,\quad
 \deg r_a=88,\quad\deg r_b=102,\quad\deg r_e=114,\quad
 \deg A_c=45.
\]
Put $z_0=Sn-A_c$, of degree 45.  For $P_R$, the relation
$C_y=(\delta+1)V^2-1$ gives the exact cancellation
\[
 \delta(SV+z_0)^2-S^2C_y
   =\delta z_0^2+2\delta SVz_0-S^2(V^2-1).
\]
The right side has degree at most 136.  Thus the norm register $q_0$ has
unique highest part $\delta_0(r_e)_0$ of degree 146, while
$(q_1)_0=-2\delta_0S_0V_0$ has degree 91.  Subsequent strict degree
comparisons give
\[
\begin{array}{c|rrrrrrrrr}
 &q_0&q_1&T_0&R_0&v_0&v_1&n_0&n_1&P_R\\\hline
 \deg&146&91&254&182&296&237&592&533&1184.
\end{array}
\]
Here and below a subscript zero on a polynomial denotes its highest
homogeneous component, independently of the temporary norm-register
names.  The full top form is
\[
 (P_R)_0=2^8 3^{24}S_0^{16}\delta_0^8V_0^8(D_g)_0^4.
\]
For $P_I$, the analogous identity is
\[
 (S\delta V+z_0)^2-S^2\delta C_y
 =z_0^2+2S\delta Vz_0-S^2\delta V^2+S^2\delta.
\]
Its unique highest term is $-S_0^2\delta_0V_0^2$, of degree 150.  The
remaining induction gives
\[
\begin{array}{c|rrrrrrrrr}
 &q_0&q_1&T_0&R_0&v_0&v_1&n_0&n_1&P_I\\\hline
 \deg&150&91&190&182&300&241&600&541&1200,
\end{array}
\]
and $(P_I)_0=S_0^{16}\delta_0^8V_0^{16}$.

These bounds are attained on the actual fixed tuple, not merely on generic
symbolic ports.  Selecting the largest $T$ power in each top form gives
\begin{align*}
 [j^{624}w^{144}v^{160}s^{64}r^{56}c^{112}h^{16}T^8]P_R
   &=2^{1256}3^{796}\beta^{624}>0,\\
 [j^{576}w^{144}v^{144}s^{80}r^{80}c^{160}T^{16}]P_I
   &=2^{1200}3^{720}\beta^{576}>0.
\end{align*}
The exponent sums are 1184 and 1200.  These are distinguished top-degree
monomials, not assertions that the full homogeneous components contain
only one term.  In particular the top part of $V$ contains
$4I_0T-R_0$, whose two summands both have degree five.  The exact checker
keeps the graded-lex leading term with $T$ first and confirms every
coefficient induction after the displayed cancellation.

\subsection{Cost and what the computations certify}

Under the same fully charged binary-operation convention as the parent,
the literal counts are
\begin{center}
\begin{tabular}{lrrrrr}
Equation&Multiplications&Additions&Subtractions&Total&Degree\\\hline
Old rational&93&26&26&145&1200\\
New $P_R$&93&27&26&146&1184\\
Old integral&89&26&26&141&1456\\
New $P_I$&90&27&26&143&1200
\end{tabular}
\end{center}
Thus $P_I$ reaches the old rational degree with two fewer operations,
while $P_R$ supplies a lower-degree option within this historical
comparison.  The saturated $123/1084$ equation of
Section~\ref{sec:integer12-saturated} improves both metrics of every row.

The standard-library verifier imports no new producer.  It proves the
whole-output substitution identities by a common expression DAG, checks
both integer inverse maps, binds all prefix formulas by independent
free-ring cuts, verifies every generic norm coefficient, and performs
the exact leading induction.  It also reconstructs all fixed coefficients
from the finite recipe and checks the live operation counts.  Normal and
optimized Python runs pass, and corrupted center and norm signs are
rejected.  These checks establish exact polynomial algebra and
transcription; the semantic conclusion uses the unbounded bijection proof
and the named ordinary parent theorem.
